\begin{onehalfspace}

\section{Dal modello ideale alla macchina fisica}

Shor è definito nel modello ideale di porte unitarie e misure perfette. I processori non fault-tolerant operano invece nel regime \ac{NISQ}~\cite{preskill2018}: coerenza finita, gate imperfetti, connettività limitata e lettura rumorosa restringono la profondità utile. Il numero di qubit, da solo, non determina quindi la capacità di eseguire un circuito.

La protezione non può consistere nella copia diretta dello stato, vietata dal teorema del no-cloning~\cite{wootters1982}. Inoltre un errore quantistico può essere una rotazione continua o un processo dissipativo, non soltanto l'inversione discreta di un bit. Per collegare fenomeno fisico e simulazione occorre distinguere la sorgente dell'errore dal \emph{canale} matematico con cui la si approssima.

\section{I vincoli hardware rilevanti}
\label{sec:anatomia_qpu}

Per il seguito sono sufficienti tre proprietà delle \ac{QPU} superconduttive. Primo, i gate sono impulsi analogici: errori di calibrazione, accoppiamenti residui e leakage perturbano l'unitaria desiderata. Secondo, ogni operazione ha una durata, mentre rilassamento e perdita di fase avvengono su scale $T_1$ e $T_2$; la durata del cammino critico deve quindi restare piccola rispetto ai tempi di coerenza. Terzo, le interazioni a due qubit rispettano una topologia fisica: il routing aggiunge operazioni e l'esecuzione parallela può generare crosstalk. Anatomia criogenica, controllo e tecnologie di qubit sono approfonditi nella Sezione~\ref{app:anatomia} dell'Appendice~\ref{app:teoria}.

\section{Fonti fisiche e canali adottati}
\label{sec:fonti_fisiche}
\label{sec:canali}

Un canale quantistico è una mappa lineare \acf{CPTP} sulla matrice densità. Nella rappresentazione di Kraus~\cite{nielsenChuang2000},
\begin{equation}
    \mathcal{E}(\rho)=\sum_k K_k\rho K_k^\dagger,
    \qquad \sum_k K_k^\dagger K_k=I.
    \label{eq:kraus}
\end{equation}
La Tabella~\ref{tab:mapping_rumore} riunisce i livelli fisico, matematico e algoritmico, evitando di trattare come equivalenti modelli che descrivono fenomeni diversi.

\begin{table}[!htbp]
\centering
\small
\begin{tabular}{@{}p{2.5cm}p{3.2cm}p{2.8cm}p{4.0cm}@{}}
\toprule
\textbf{Fonte} & \textbf{Modello} & \textbf{Parametro} & \textbf{Impatto principale su Shor} \\
\midrule
Rilassamento energetico
& amplitude damping / thermal relaxation
& $\gamma=1-e^{-t/T_1}$
& perdita della popolazione $\ket{1}$ durante i blocchi profondi \\[2pt]
Perdita di fase
& phase flip / thermal relaxation
& $p=\tfrac12(1-e^{-t/T_2})$
& attenuazione delle coerenze usate dalla QPE \\[2pt]
Imprecisione dei gate
& canale depolarizzante a uno o due qubit
& $\lambda_{1q},\lambda_{2q}$
& accumulo sulle porte entangling e distorsione delle fasi \\[2pt]
Errore di lettura
& matrice di confusione classica
& $e_0,e_1$
& alterazione dei bit da cui le frazioni continue stimano $r$ \\[2pt]
Crosstalk
& errore correlato o coerente
& dipendente da coppia e calendario
& correlazioni non descritte dal modello uniforme indipendente \\
\bottomrule
\end{tabular}
\caption{Corrispondenza fra principali fonti di rumore, modelli e vulnerabilità dell'algoritmo.}
\label{tab:fonti_rumore}
\label{tab:mapping_rumore}
\label{tab:canali_rumore}
\end{table}

\subsection{Rilassamento e perdita di fase}

Il tempo $T_1$ descrive il decadimento $\ket{1}\rightarrow\ket{0}$, mentre $T_2$ descrive la perdita della fase relativa. Separando la componente puramente dissipativa da quella di puro dephasing,
\begin{equation}
    \frac{1}{T_2}=\frac{1}{2T_1}+\frac{1}{T_\phi},
    \label{eq:T2_decomp}
\end{equation}
da cui $T_2\leq2T_1$. Il canale di amplitude damping è non unital: sposta lo stato verso $\ket{0}$. Il phase flip lascia invece invariate le popolazioni e moltiplica le coerenze per $1-2p$. Nella simulazione, Qiskit Aer combina i due contributi nel \emph{thermal relaxation error}, parametrizzato da $T_1$, $T_2$ e dalla durata del gate. Il modello è Markoviano e non include automaticamente deriva, leakage o correlazioni temporali.

\subsection{Canale depolarizzante}

La campagna usa il canale depolarizzante per rappresentare l'errore stocastico associato ai gate. Per un qubit, indicando con $p$ la probabilità totale di un Pauli non identità,
\begin{equation}
    \mathcal{D}_{p}(\rho)
    =(1-p)\rho+\frac{p}{3}
      \left(X\rho X+Y\rho Y+Z\rho Z\right).
    \label{eq:depolarizing}
\end{equation}
Il canale contrae isotropicamente il vettore di Bloch e non descrive la direzione microscopica dell'errore. Aer usa invece il parametro $\lambda$ nella convenzione
\begin{equation}
    \mathcal{D}_{\lambda}^{(q)}(\rho)
    =(1-\lambda)\rho+\lambda\,\frac{I}{2^q},
    \label{eq:depolarizing_compact}
\end{equation}
per un gate su $q$ qubit. La probabilità aggregata della componente Pauli non identità è dunque $3\lambda/4$ per $q=1$ e $15\lambda/16$ per $q=2$. Esplicitare questa convenzione è necessario per interpretare correttamente i parametri sperimentali; la definizione del noise model e i gate ai quali è applicato sono riportati nel Capitolo~\ref{chap:metodologia}.

\subsection{Lettura e rumore correlato}

L'errore di misura è rappresentato, per ciascun qubit, dalla matrice di assegnazione
\begin{equation}
    A=
    \begin{pmatrix}
        1-e_0 & e_1\\
        e_0   & 1-e_1
    \end{pmatrix},
    \qquad
    A_{ij}=P(\text{esito }i\mid\text{stato }j).
\end{equation}
Per errori indipendenti su più qubit si usa il prodotto tensoriale delle matrici locali. Il crosstalk viola invece l'ipotesi di indipendenza: un'operazione può perturbare qubit vicini o dipendere dalle operazioni simultanee~\cite{sarovar2020}. Il modello uniforme della prima campagna non pretende di riprodurre tali correlazioni; queste vengono introdotte separatamente negli esperimenti di decodifica, dove è possibile confrontare metodi analitici e appresi sullo stesso insieme di sindromi.

Le derivazioni degli operatori di Kraus, l'azione sulle componenti della matrice densità e la geometria sulla sfera di Bloch sono raccolte nella Sezione~\ref{app:canali} dell'Appendice~\ref{app:teoria}.

\section{Perché Shor è sensibile al rumore}

La vulnerabilità non deriva da un solo canale, ma dall'interazione fra profondità, precisione di fase e misura:
\begin{itemize}
    \item l'\textbf{esponenziazione modulare} domina la profondità e contiene molte porte a due qubit; aumentano sia l'esposizione agli errori di gate sia il tempo durante il quale i registri devono restare coerenti;
    \item la \textbf{phase estimation} codifica il periodo nelle fasi relative. Dephasing ed errori coerenti riducono o spostano i picchi da cui si ricava $s/r$;
    \item la \textbf{QFT inversa} contiene rotazioni controllate, comprese rotazioni ad angolo piccolo, sensibili alla precisione di controllo;
    \item la \textbf{misura} finale può modificare i bit dell'intero osservato e quindi la frazione continua selezionata;
    \item \textbf{routing e parallelismo} aggiungono gate e possono introdurre errori correlati dipendenti dalla topologia.
\end{itemize}

\begin{figure}[!htbp]
\centering
\resizebox{0.95\textwidth}{!}{%
\begin{tikzpicture}[
  >=stealth,
  stage/.style={rectangle,draw,rounded corners=2pt,fill=blue!10,
                minimum height=12mm,text width=28mm,align=center,font=\small},
  flow/.style={->,thick}
]
\node[stage] (S1) at (0,0)   {Preparazione\\$H^{\otimes t}$};
\node[stage] (S2) at (4,0)   {Esponenziazione\\modulare};
\node[stage] (S3) at (8,0)   {\ac{QFT}$^{-1}$};
\node[stage] (S4) at (12,0)  {Misura e\\post-processing};
\draw[flow] (S1)--(S2);
\draw[flow] (S2)--(S3);
\draw[flow] (S3)--(S4);
\node[below=7mm of S2,align=center,font=\scriptsize,text=red!60!black]
  {decoerenza, gate 2Q,\\crosstalk};
\node[below=7mm of S3,align=center,font=\scriptsize,text=violet!70!black]
  {dephasing, errori\\di rotazione};
\node[below=7mm of S4,align=center,font=\scriptsize,text=orange!80!black]
  {readout};
\end{tikzpicture}%
}
\caption[Localizzazione del rumore nella pipeline di Shor]{Principali sorgenti di errore lungo la pipeline; decoerenza e crosstalk possono agire trasversalmente, non soltanto nello stadio indicato.}
\label{fig:errori_shor}
\end{figure}

Non è corretto trasformare il solo conteggio asintotico $O((\log N)^3)$ in una probabilità di successo: gate diversi hanno durate e tassi diversi, il circuito contiene parallelismo e alcuni errori non cambiano l'esito classico. Per questa ragione il seguito usa circuiti compilati e noise model dichiarati, separando conteggio di gate, profondità, parametri fisici e metrica di successo. Le dimostrazioni di fattorizzazione su istanze piccole e circuiti semplificati~\cite{skosana2021} non rimuovono questo limite di scalabilità.

Il capitolo successivo distingue quindi gli interventi possibili: riduzione fisica del rumore, mitigazione a posteriori, correzione/tolleranza ai guasti e metodi appresi. Nessuno di essi viene assunto efficace per definizione; ciascuno richiede un confronto entro il proprio livello operativo.

\end{onehalfspace}
